% Generated draft: requires mathematical and source review.
\documentclass[11pt]{article}
\usepackage[margin=25mm]{geometry}
\usepackage{fontspec}
\setmainfont{Latin Modern Roman}
\usepackage{amsmath,amssymb,mathtools,unicode-math}
\setmathfont{Latin Modern Math}
\usepackage{enumitem,needspace,xurl,hyperref}
\setlength{\parindent}{0pt}
\setlength{\parskip}{5pt}
\setlength{\emergencystretch}{3em}
\newcommand{\N}{\mathbb{N}}
\newcommand{\Z}{\mathbb{Z}}
\newcommand{\Q}{\mathbb{Q}}
\newcommand{\R}{\mathbb{R}}
\newcommand{\C}{\mathbb{C}}
\newcommand{\F}{\mathbb{F}}
\newcommand{\A}{\mathbb{A}}
\newcommand{\PP}{\mathbb{P}}
\newcommand{\E}{\mathbb{E}}
\newcommand{\eps}{\varepsilon}
\newcommand{\dd}{\,\mathrm{d}}
\begin{document}
\section*{UnsolvedMath ID: 20001956}
\textbf{Source problem code:} catalogue entry
\subsection{Definitions and mathematical statement}
E. Lerman's Question 5.1 asks for a hyperkähler analogue of the topological mechanisms behind Kirwan surjectivity, allowing maps more general than geometric moment maps. The source invokes [LT97] for a small-resolution calculation using a map that is not a moment map, and [GGK02] for abstract moment maps on \(G\)-manifolds without assuming a symplectic form. An abstract moment map \(\mu:M\to\mathfrak g^*\) is described as equivariant and subject to subgroup/fixed-point conditions that are not reproduced in the record. Thus no complete definition or single fully specified general conjecture is supplied.

The excerpt proposes componentwise nondegeneracy: component functions are Morse--Bott and their critical sets correspond to fixed-point sets. It writes \(M^G\) for these sets, without stating which components or subgroup conventions are intended. For a component \(\mu^\xi=\langle\mu,\xi\rangle\), full-group fixed points and fixed points of the subgroup generated by \(\xi\) need not agree; the zero component also needs exclusion in a nontrivial action. These qualifications cannot be silently supplied as the exact [GGK02] definition. The rotating \(\R^2\) with a constant map is the source's example showing why the abstract axioms alone do not ensure the desired nondegeneracy condition.

The supplied aid gives the following conditional topological certificate. Let \(G\) be a compact Lie group, \(\Bbbk\) a field, and \(\mu\) an equivariant map with zero set \(Z\). Assume a finite \(G\)-equivariant filtration
\[
 Z=X_0\subset X_1\subset\cdots\subset X_N=M
\]
whose successive stages are Thom attachments over \(G\)-spaces \(C_j\), with negative bundles \(\nu_j\) of rank \(\lambda_j\). Require the Thom/excision identification for each pair, with the restriction to the zero section identified with multiplication by the equivariant Euler class. Assume \(\nu_j\) is equivariantly \(\Bbbk\)-oriented and
\[
 e_G(\nu_j):H_G^{q-\lambda_j}(C_j;\Bbbk)\longrightarrow H_G^q(C_j;\Bbbk)
\]
is injective for every \(q\). Here \(H_G^*(X;\Bbbk)=H^*(EG\times_GX;\Bbbk)\). The reported stagewise conclusion is
\[
 0\longrightarrow H_G^{q-\lambda_j}(C_j;\Bbbk)
 \longrightarrow H_G^q(X_j;\Bbbk)
 \longrightarrow H_G^q(X_{j-1};\Bbbk)\longrightarrow0,
\]
and hence surjectivity of \(H_G^*(M;\Bbbk)\to H_G^*(Z;\Bbbk)\).

For a hyperkähler moment map \(\boldsymbol\mu=(\mu_I,\mu_J,\mu_K):M\to\mathfrak g^*\otimes\R^3\), choose an invariant norm and consider \(F=\|\boldsymbol\mu\|^2\), with \(Z=\boldsymbol\mu^{-1}(0)\). The aid states the certificate as a conditional Kirwan-type criterion when the required finite filtration comes from \(F\). Its two layers are distinct: construct the filtration with its Thom properties, then check orientations and Euler-class injectivity over the chosen field. It does not assert these hypotheses for every hyperkähler or abstract moment map. Identifying Borel cohomology of \(Z\) with ordinary cohomology of a geometric quotient requires additional action hypotheses not given here.

The aid's obstruction to a simpler criterion is a shifted-height map on the rotating sphere. An explicit representative is \(S^2\subset\R^3\), with \(S^1\) rotating about the vertical axis, height \(h(x,y,z)=z\), and \(\mu=h-1\). Every nonzero scalar component of this one-dimensional map is Morse with critical set equal to the two poles. Nevertheless \(F=(h-1)^2\) has an isolated critical minimum at the north pole with vanishing Hessian and is not Morse--Bott. Thus componentwise nondegeneracy does not supply the missing norm-square or filtration hypothesis. This illustrates the reported obstruction, not a counterexample to surjectivity under the certificate's full assumptions.
\subsection{Short English statement}
The source asks which topological features of moment maps can yield hyperkähler Kirwan-type surjectivity. The aid separates a finite Thom-filtration hypothesis from coefficient-sensitive orientation and Euler-class checks. A shifted height function shows why nondegeneracy of individual components cannot replace a hypothesis on the norm-square.
\subsection{Sources}
\begin{itemize}
\item Local UnsolvedMath registry, permanent ID 20001956, exact source code catalogue entry; canonical \texttt{aim-geometry-notes.json}, \texttt{notes[293]}, tag \texttt{question}; AIM Workshop Problem Lists, Geometry, source item 5.1, attributed to E. Lerman. Created and updated 14 August 2026.
\item AIM list \emph{Moment maps and surjectivity in various geometries}, Question 5.1; \url{https://aimath.org/WWN/momentmaps/momentmaps.pdf}. The supplied verification identifies this attribution and location; the PDF was not accessed.
\item Original reference labels [LT97], for topology of a small resolution of a singular symplectic reduced space using a non-moment map, and [GGK02], for abstract moment maps and nondegeneracy. Complete bibliographic entries and the full abstract axioms are not supplied.
\item Supplied attempt 1: \texttt{partial\_result}, problem status at run \texttt{partially\_solved}, rigor label \texttt{proved}, novelty confidence \texttt{low}; classification \texttt{PARTIAL-PROGRESS}, registry status \texttt{partially\_solved}. Candidate: finite-filtration certificate and componentwise-nondegeneracy obstruction. These are source labels, not independent certification.
\item Ulam AI / UnsolvedMath Contributors; CC BY 4.0. Machine-generated research aid, not peer reviewed; independently verify before citation.
\end{itemize}
\end{document}
